%% ELEC2013 2007

%% ----------------------------------------------------------------
%% exam.tex
%% ----------------------------------------------------------------
\documentclass{sotonExamBoxes}    % use option [answers] to include them
%\documentclass[answers]{uosexamb}
%% ----------------------------------------------------------------
\usepackage{graphicx}
%\usepackage{pstricks}
\usepackage{alltt}
\usepackage{bm}
\usepackage{amssymb}
\usepackage{amsfonts}
\usepackage{amsmath}
\thinmuskip=1mu
%\renewcommand{\ttdefault}{pcr}
%% macros.tex -- shared macro definitions for the problem sheets.
%%
%% Every sheet inputs this instead of carrying its own copy. Input it after
%% the \usepackage lines: the definitions below use bm, amsmath and amssymb.
%%
%% Six macros were defined inconsistently across the sheets when this file
%% was made; the choices are marked CONFLICT below.




\usepackage{stmaryrd}

\renewcommand{\ttdefault}{pcr}

\newcommand{\tr}{\textsf{T}}

\newcommand{\e}[1]{{\rm e}^{#1}}

\newcommand{\E}[1]{{\rm e}^{{\displaystyle #1}}}

\newcommand{\logg}[1]{\log\left(#1\right)}

\newcommand{\Prob}[1]{\mathbb{P}\left(#1\right)}

\newcommand{\Step}{\mathop{\mathrm{O}\hspace{-15pt}\rule[8.5pt]{11pt}{1.2pt}\hspace{2pt}}}

\newcommand\diag{\mathop{\mathrm{diag}}}

\newcommand{\grad}{\bm{\nabla}}

\newcommand{\dd}{\mathrm{d}}

\newcommand{\class}{\mathcal{C}}

\newcommand{\data}{\mathcal{D}}

%% CONFLICT: gaussianProcess.tex had {U}{eur}{b}{n} (Euler bold);
%% every other sheet had this. Matrices in that one sheet now look
%% different -- see the report.

\DeclareMathAlphabet{\mat}{OT1}{cmss}{bx}{n}

%% The identity needs to be distinguishable. In cmss a bold I is a plain
%% vertical bar with no crosspieces, so \mat{I} reads as a rule rather than a
%% letter. Euler's bold I keeps its serifs, so the identity alone is set in
%% Euler while everything else stays cmss.
\DeclareMathAlphabet{\matEuler}{U}{eur}{b}{n}
\newcommand{\Id}{\matEuler{I}}

%% CONFLICT and BUG FIX: 16 sheets defined this using \qpartial,
%% which is defined nowhere, so \pd was broken in all of them.
%% Only problem2.tex had the working \partial form, kept here.

\newcommand{\pd}[2]{\ensuremath{\frac{\partial #1}{\partial #2}}\xspace}

\newcommand{\M}[1]{\ensuremath{\boldsymbol{#1}}\xspace}

\newcommand{\bst}{\rule{0mm}{5mm}}

%% CONFLICT: some sheets omitted the \! . Kept the tighter spacing,
%% which is what every sheet that actually uses \av had.

\newcommand{\av}[2][\,]{\mathbb{E}_{#1}\!\left[ {\strut #2} \right]}

\newcommand{\len}[1]{\| #1 \|}

\newcommand{\inner}[2]{\left\langle #1, #2\right\rangle}

\newcommand{\squeeze}{\setlength{\itemsep}{-2pt}}

%% CONFLICT: two spacings, identical output in maths mode.

\newcommand{\bra}[1]{\left( #1 \right)}

\newcommand{\pred}[1]{\left\llbracket { \small #1} \right\rrbracket}

\newcommand{\mln}{\hat{f}_n(\bm{x})}


%% \mathop makes "tr" an operator rather than three variables juxtaposed:
%% it is set upright and TeX adds its own spacing on either side. \nolimits
%% keeps a sub/superscript beside it rather than under it in a displayed
%% equation.
%%
%% The space is \; (thickmuskip) and not \, (thinmuskip) because every sheet
%% sets \thinmuskip=1mu, which would leave \, all but invisible. Use \: for a
%% slightly tighter gap.
\newcommand{\Tr}{\mathop{\mathrm{tr}}\nolimits\;}

%% NOT SHARED: gaussianProcess.tex redefines \det to take an argument
%% (\det{\mat{V}}), while calculus.tex and vectorSpaces.tex use the standard
%% no-argument operator (\det(\mat{X})). The two cannot coexist, so that
%% redefinition stays local to gaussianProcess.tex.

%% NOT SHARED: \ml and \mm mean different things in different sheets.
%% amgm and convexity use a no-argument \ml and \mm = \hat{m}_c(x);
%% ensembleLearning uses a one-argument \ml and \mm = \hat{m}(x).
%% Both forms are in use, so each sheet keeps its own.

\graphicspath{{figures/} {../2018-19/figures}}
\graphicspath{{figures/}}

% \usepackage{picture}
\begin{document}
\unitTitle{Vector Space Problem Sheet}
\authors{Adam Pr\"ugel-Bennett}
\unitcode{COMP6208}
\semester{2}
\year{2022}
\duration{45}{0}


\maketitle

This paper asks you to prove some well known results.  Although the
algebra is easy the proofs are not entirely straightforward.  There
are marks assigned to the readability of the solution and also how
well laid out and explained the steps you make are. (A good proof
needs to be easy to follow: you need not comment on trivial algebra,
but there should not be steps that are difficult to follow).

This looks very mathematical, but it helps to develop the tools and
language that is used to describe machine learning.

%\blankpage


%% ----------------------------------------------------------------
%\section
% -----------------------------------------------------------------------
% *************************  Section A  *********************************
% -----------------------------------------------------------------------

% ***************************** question 1 **************************


\begin{question}{20}
  An inner product $\inner{\bm{x}}{\bm{y}}$ between vectors in a
  vector space $\mathcal{V}$ satisfies the following
  properties
  \begin{enumerate}\squeeze
  \item $\inner{\bm{x}}{\bm{x}} \geq 0$ for all $\bm{x}\in\mathcal{V}$
  \item $\inner{\bm{x}}{\bm{x}} = 0$  if and only if $\bm{x}=\bm{0}$
  \item $\inner{\alpha\, \bm{x}}{\bm{y}} = \alpha\,
    \inner{\bm{x}}{\bm{y}}$
  \item $\inner{\bm{x}}{\bm{y}+\bm{z}} = \inner{\bm{x}}{\bm{y}} +
    \inner{\bm{x}}{\bm{z}}$
  \item $\inner{\bm{x}}{\bm{y}} = \inner{\bm{y}}{\bm{x}}$
  \end{enumerate}
  The question explores properties of inner products.
  \clearpage
  \begin{qparts}
    \qpart[10]{Consider the quadratic
      \begin{align*}
        q(t ) = \inner{\bm{x} + t\, \bm{y}}{\bm{x} + t\, \bm{y}}
      \end{align*}
      By definition 1 of an inner product $q(t)$ must be non-negative
      and will only be 0 when $\bm{x} + t\, \bm{y}=0$.

      Expand $q(t)$ in the form $q(t) = A\,t^2 + 2\, B \, t +C$ to       
      find $A$, $B$ and $C$.  For $q(t)$ to not change sign its roots
      (values of $t$ such that $q(t)=0$) must be complex (i.e. have an
      imaginary part), or possible have a double root if there is a
      value of $t$ such that $q(t)=0$.  Use the
      standard solutions to a quadratic of the form $q(t) = A\,t^2 +
      2\, B t + C$ to show that for our $q(t)$ to never become
      negative implies
      \begin{align*}
        \inner{\bm{x}}{\bm{y}}^2 \leq \inner{\bm{x}}{\bm{x}} \,
        \inner{\bm{y}}{\bm{y}}.
      \end{align*}
      This is the famous Cauchy-Schwarz inequality written in a very
      general form.
      \begin{answer}    
        \begin{align*}
          q(t) = \inner{\bm{y}}{\bm{y}}  t^2 + 2\,
          \inner{\bm{x}}{\bm{y}}\,t +  \inner{\bm{x}}{\bm{x}}
        \end{align*}
        So that $A= \inner{\bm{y}}{\bm{y}}$, $B= 2\,\inner{\bm{x}}{\bm{y}}$
        and $C=\inner{\bm{x}}{\bm{x}}$.  The roots to the quadratic
        equation $q(t)=A\,t^2 + B \, t + C = 0$ are
        \begin{align*}
          t = \frac{-B \pm \sqrt{B^2 - 4\,A\,C}}{2\,A}.
        \end{align*}
        The condition that these roots are complex is that
        $B^2-4\,A\,C<0$, and where $B^2-4\,A\,C=0$ indicating there is a
        double root.  Thus $B^2/4 \leq A\,C$.  Putting in the values of
        $A$, $B$ and $C$ then
        \begin{align*}
          \inner{\bm{x}}{\bm{y}}^2 \leq \inner{\bm{x}}{\bm{x}} \,
          \inner{\bm{y}}{\bm{y}}.
        \end{align*}
      \end{answer}}{\lines{17}}
  
    \qpart[10]{From an inner product we can define a norm $\len{\bm{x}} =
      \sqrt{\inner{\bm{x}}{\bm{x}}}$.  This clearly satisfies
      non-negativity and linearity.  The only non-trivial property to
      show is that this norm satisfies the triangular inequality.  Expand out
      $\len{\bm{x} + \bm{y}}$ and hence show that
      \begin{align*}
        \len{\bm{x}+\bm{y}}^2 = \len{\bm{x}}^2 + \len{\bm{y}}^2 + 2
        \inner{\bm{x}}{\bm{y}}
      \end{align*}
      Then use the Cauchy-Schwarz inequality to prove the triangular
      inequality.
    \begin{answer}
      We start form
      \begin{align*}
        \len{\bm{x}+\bm{y}}^2
        &=  \inner{\bm{x} + \bm{y}}{\bm{x} + \bm{y}}
        \\
        &=  \inner{\bm{x}}{\bm{x}} + \inner{\bm{y}}{\bm{y}}
          + 2 \, \inner{\bm{x}}{\bm{y}}
      \end{align*}
      Using the Cauchy-Schwarz inequality $|\inner{\bm{x}}{\bm{y}}|^2
      \leq \inner{\bm{x}}{\bm{x}}\, \inner{\bm{y}}{\bm{y}}$ so that
      \begin{align*}
         \len{\bm{x}+\bm{y}}^2
        &\leq  \inner{\bm{x}}{\bm{x}} + \inner{\bm{y}}{\bm{y}}
          + 2 \, \sqrt{\inner{\bm{x}}{\bm{x}}\, \inner{\bm{y}}{\bm{y}}}
        \\
        &= \len{\bm{x}}^2 + \len{\bm{y}}^2 + 2 \,
          \len{\bm{x}}\,\len{\bm{y}}
        = \left( \len{\bm{x}} + \len{\bm{y}}\right)^2
      \end{align*}
      Taking the square root (which is monotonic so does not change
      the inequality)
      \begin{align*}
         \len{\bm{x}+\bm{y}} \leq  \len{\bm{x}} + \len{\bm{y}}
      \end{align*}
      Thus proving the triangular inequality for
      $\len{\bm{x}} = \sqrt{\inner{\bm{x}}{\bm{x}}}$. 
    \end{answer}}{\lines{17}}
  \end{qparts}
\end{question}
\freshpage

% ***************************** question 2 **************************

\begin{question}{20}
  Random variables, $X$, $Y$, etc.{} form a vector space (i.e. they
  satisfy properties such as closure under addition and scalar
  multiplication).   Furthermore we can define an inner product
  between random variables as
  \begin{align*}
    \inner{X}{Y} = \av{X\,Y}
  \end{align*}
  (i.e. the expectation of the random variable $Z = X\,Y$)
  \begin{qparts}
    \qpart[10]{Use the Cauchy-Schwarz inequality to show that
      \begin{align*}
        \mathrm{Cov}(X,Y)^2 \leq \mathrm{Var}(X) \, \mathrm{Var}(Y)
      \end{align*}
      hence show that the Pearson correlation is between -1 and 1.
    \begin{answer}
      Let $\mu = \av{X}$ and $\nu = \av{Y}$ then
      \begin{align*}
        \mathrm{Cov}(X,Y) = \av{(X-\mu)(Y-\nu)}
      \end{align*}
      From the Cauchy-Schwarz inequality
      \begin{align*}
        \mathrm{Cov}(X,Y)^2
        &= \av{(X-\mu)\,(Y-\nu)}^2 = \inner{X-\mu}{Y-\mu}^2
        \\
        &\leq \inner{X-\mu}{X-\mu} \, \inner{Y-\mu}{Y-\mu} =
          \av{(X-\mu)^2} \, \av{(Y-\nu)^2} = \mathrm{Var}(X)\,
          \mathrm{Var}(Y).
      \end{align*}
      which proves the inequality.  The Pearson correlation is defined
      as
      \begin{align*}
        \rho(X,Y) =  \frac{\mathrm{Cov}(X,Y)}
        {\sqrt{\mathrm{Var}(X) \, \mathrm{Var}(Y)}}
      \end{align*}
      Thus $\rho(X,Y)^2= \mathrm{Cov}(X,Y) ^2/\left( \mathrm{Var}(X) \,
        \mathrm{Var}(Y)\right)$, but  we have already shown that
      $\mathrm{Cov}(X,Y)^2 \leq \mathrm{Var}(X) \, \mathrm{Var}(Y)$ so
      $\rho(X,Y)^2\leq 1$ and $\rho(X,Y)$ will lie in the interval $[-1,1]$.
    \end{answer}}{\lines{17}}

    \qpart[10]{Show that for vectors $\bm{x}, \bm{y}\in\mathbb{R}^n$ that the
      inner product
      \begin{align*}
        \inner{\bm{x}}{\bm{y}} = \sum_{i=1}^n w_i \, x_i \,y_i
      \end{align*}
      satisfies the condition of an inner product provided $w_i>0$ for
      all $i$.  Write down the norm and distance induced by this
      inner product and provide an interpretation of what this
      distance means.
    \begin{answer}
      We consider the five condition on the inner product
      \begin{enumerate}
      \item Non-negativity follows from
        \begin{align*}
          \inner{\bm{x}}{\bm{x}} = \sum_{i=1}^n w_i \, x_i ^2
          \geq 0
        \end{align*}
        since for each term in the sum $w_i \, x_i ^2\geq0$.
      \item For $ \inner{\bm{x}}{\bm{x}} = 0$ each term in the sum
        must be zero.  Since $w_i>0$ this implies $x_i=0$ for all $i$,
        that is, $\bm{x}=\bm{0}$.
      \item Linearity follows trivially since
        \begin{align*}
          \inner{a\,\bm{x}}{\bm{y}}
          &= \sum_{i=1}^n w_i \, (a\,x_i)  \,y_i 
          & a\, \sum_{i=1}^n w_i \, x_i \,y_i = a\,\inner{\bm{x}}{\bm{y}}
        \end{align*}
      \item The distributive law again follows trivially
        \begin{align*}
          \inner{\bm{x}}{\bm{y}+\bm{z}}
          &= \sum_{i=1}^n w_i \, x_i \,(y_i+z_i) \\
          &= \sum_{i=1}^n w_i \, x_i \,y_i + \sum_{i=1}^n w_i \, x_i \,z_i 
          =   \inner{\bm{x}}{\bm{y}} +   \inner{\bm{x}}{\bm{z}}
        \end{align*}
      \item Finally symmetry follows from the commutativity of
        multiplication
        \begin{align*}
          \inner{\bm{x}}{\bm{y}} &=  \sum_{i=1}^n w_i \, x_i \,y_i
          =  \sum_{i=1}^n w_i \, y_i \,x_i =  \inner{\bm{y}}{\bm{x}} 
        \end{align*}
        The induced norm and distance for this inner product is
        \begin{align*}
          \len{x} &= \sqrt{\sum_{i=1}^n w_i \, x_i^2}, &
          d(\bm{x}, \bm{y}) &= \sqrt{ \sum_{i=1}^n w_i \, (x_i-y_i)^2}.
        \end{align*}
        A natural interpretation of this is that this is a
        generalisation of the Euclidean distance where we rescale each
        axis by a factor $w_i$ and then compute the normal Euclidean distance.
      \end{enumerate}
    \end{answer}}{\lines{19}}
  \end{qparts}
  
\end{question}




\freshpage

% ***************************** question 3 **************************


\begin{question}{20}
  In the previous question you showed that $\inner{\bm{x}}{\bm{y}} =
  \sum_{i=1}^n w_i\, x_i\, y_i$ is an inner product provided every
  weight $w_i$ is positive.  Functions form a vector space in just the
  same way (we can add them and multiply them by scalars), and if we
  replace the sum over components by an integral over $x$ we get
  \begin{align*}
    \inner{f}{g} = \int_{-\infty}^{\infty} w(x)\, f(x)\, g(x)\, \dd x
  \end{align*}
  which, by the same argument, is an inner product provided the
  \emph{weight function} $w(x)$ is positive.

  In this question we choose $w(x)$ to be the density of a standard
  normal deviate
  \begin{align*}
    w(x) = \frac{1}{\sqrt{2\,\pi}}\, \e{-x^2/2}
  \end{align*}
  so that $\inner{f}{g} = \av{f(X)\, g(X)}$ where $X\sim
  \mathcal{N}(0,1)$.  That is, this is exactly the inner product of
  the previous question, applied to functions of a random variable.

  The monomials $f_i(x) = x^i$ span the polynomials, but they are not
  orthogonal.  We can orthogonalise them with the Gram--Schmidt
  procedure: having found $p_0, p_1, \ldots, p_{n-1}$ we subtract from
  $f_n$ its component along each of them,
  \begin{align*}
    p_n(x) = f_n(x) - \sum_{i=0}^{n-1}
    \frac{\inner{f_n}{p_i}}{\inner{p_i}{p_i}}\, p_i(x) ,
    \qquad p_0(x) = 1 .
  \end{align*}

  Orthogonal polynomials are useful because they make approximating a
  function easy.  If we want the best polynomial approximation to some
  function $f$ we can compute each coefficient separately with a
  single inner product, and if we later decide to add a higher order
  term none of the coefficients we have already found change.  Using
  the monomials instead we would have to solve a set of simultaneous
  equations, and start again from scratch each time we added a term.

  You will need the integrals below, which you are not expected to
  compute.  (They are the moments of a standard normal deviate.  All
  the odd moments vanish because $w(x)\,x^n$ is an odd function.)
  \begin{align*}
    \begin{array}{c|ccccccc}
      n & 0 & 1 & 2 & 3 & 4 & 5 & 6 \\ \hline
      \av{X^n} = \int_{-\infty}^{\infty} w(x)\, x^n\, \dd x
        & 1 & 0 & 1 & 0 & 3 & 0 & 15
    \end{array}
  \end{align*}

  \begin{qparts}

    \qpart[5]{Show that $\inner{f}{g} = \int w(x)\, f(x)\, g(x)\, \dd
      x$ satisfies the five conditions for an inner product given at
      the start of this sheet, provided $w(x)>0$.
      \begin{answer}
        The argument is the same as for the weighted sum, with the
        sum over $i$ replaced by the integral over $x$.
        \begin{enumerate}\squeeze
        \item $\inner{f}{f} = \int w(x)\, f^2(x)\, \dd x \geq 0$,
          since the integrand $w(x)\, f^2(x)\geq0$ everywhere.
        \item If $\inner{f}{f} = 0$ then, as the integrand is
          non-negative and $w(x)>0$, we need $f(x)=0$ (for all but a
          set of points of zero total length).
        \item $\inner{a\,f}{g} = \int w(x)\, a\, f(x)\, g(x)\, \dd x =
          a \inner{f}{g}$.
        \item $\inner{f}{g+h} = \int w(x) f(x) \bra{g(x)+h(x)} \dd x =
          \inner{f}{g} + \inner{f}{h}$, as the integral of a sum is
          the sum of the integrals.
        \item $\inner{f}{g} = \inner{g}{f}$ since $f(x)\,g(x) =
          g(x)\,f(x)$.
        \end{enumerate}
      \end{answer}}{\lines{12}}

    \qpart[5]{Starting from $p_0(x) = 1$, use Gram--Schmidt to find
      $p_1(x)$ and $p_2(x)$.
      \begin{answer}
        For $p_1$ we subtract from $f_1(x)=x$ its component along
        $p_0$.  Since $\inner{f_1}{p_0} = \av{X} = 0$,
        \begin{align*}
          p_1(x) = x - \frac{0}{1}\, 1 = x .
        \end{align*}
        For $p_2$ we start from $f_2(x) = x^2$ and use
        $\inner{p_0}{p_0} = \av{1} = 1$ and $\inner{p_1}{p_1} =
        \av{X^2} = 1$,
        \begin{align*}
          \inner{f_2}{p_0} &= \av{X^2} = 1 , &
          \inner{f_2}{p_1} &= \av{X^3} = 0 ,
        \end{align*}
        so that
        \begin{align*}
          p_2(x) = x^2 - \frac{1}{1}\, 1 - \frac{0}{1}\, x = x^2 - 1 .
        \end{align*}
      \end{answer}}{\lines{12}}

    \qpart[5]{Find $p_3(x)$.
      \begin{answer}
        We need $\inner{p_2}{p_2} = \av{(X^2-1)^2} = \av{X^4} -
        2\av{X^2} + 1 = 3-2+1 = 2$, and starting from $f_3(x) = x^3$
        \begin{align*}
          \inner{f_3}{p_0} &= \av{X^3} = 0 , &
          \inner{f_3}{p_1} &= \av{X^4} = 3 , \\
          \inner{f_3}{p_2} &= \av{X^3(X^2-1)} = \av{X^5} - \av{X^3} = 0 .
        \end{align*}
        Therefore
        \begin{align*}
          p_3(x) = x^3 - \frac{0}{1}\,1 - \frac{3}{1}\, x -
          \frac{0}{2}\, (x^2-1) = x^3 - 3\, x .
        \end{align*}
        (These are the Hermite polynomials.  As a check, every
        $p_n$ is either an odd or an even function, which had to
        happen because $w(x)$ is symmetric about the origin.)
      \end{answer}}{\lines{12}}

    \qpart[5]{If we write a function as $f(x) = \sum_i a_i\, p_i(x)$,
      show that $a_i = \inner{f}{p_i}/\inner{p_i}{p_i}$.  Use this to
      write $f(x) = x^3$ in terms of $p_0$, $p_1$, $p_2$ and $p_3$,
      and check your answer.
      \begin{answer}
        Taking the inner product of $f = \sum_j a_j\, p_j$ with $p_i$
        \begin{align*}
          \inner{f}{p_i} = \sum_j a_j \inner{p_j}{p_i}
          = a_i \inner{p_i}{p_i}
        \end{align*}
        because $\inner{p_j}{p_i} = 0$ unless $i=j$.  Rearranging
        gives $a_i = \inner{f}{p_i}/\inner{p_i}{p_i}$.  This is the
        point of orthogonality: one integral gives one coefficient.

        For $f(x) = x^3$ we computed the three inner products we need
        in the previous part, and $\inner{f}{p_3} = \av{X^3(X^3-3X)} =
        \av{X^6} - 3\av{X^4} = 15-9 = 6$ with $\inner{p_3}{p_3} =
        \av{(X^3-3X)^2} = \av{X^6} - 6 \av{X^4} + 9 \av{X^2} =
        15-18+9=6$.  Hence $a_0 = 0$, $a_1 = 3/1 = 3$, $a_2 = 0$ and
        $a_3 = 6/6 = 1$, giving
        \begin{align*}
          x^3 = 3\, p_1(x) + p_3(x) .
        \end{align*}
        Checking, $3\,x + (x^3-3\,x) = x^3$ as required.
      \end{answer}}{\lines{14}}

  \end{qparts}
\end{question}




% ***************************** End of exam *****************************
\end{document}




